% ID: SEL-D0049-OA-E02
% ESTADO: validada. S050.
\textbf{a)} $f(4)=3$ y el límite derecho es $-8+11=3$. La rama izquierda es continua hasta $4$, así que $f$ es continua en la unión. Sus derivadas laterales son $-2\cdot4+6=-2$ y $-2$. Coinciden, por lo que $f'(4)=-2$.

\textbf{b)} La primera rama es $4-(x-3)^2$ y la segunda es un segmento decreciente. La función crece de $2$ a $3$ y decrece de $3$ a $5$. El máximo absoluto es $f(3)=4$. Como $f(2)=3$ y $f(5)=1$, el mínimo absoluto es $f(5)=1$.
\begin{center}
\setlength{\unitlength}{0.8mm}
\begin{picture}(123,58)(-5,-6)
\qbezier(0,0)(55,0)(110,0)
\qbezier(0,0)(0,24)(0,48)
\qbezier(40,30)(60,50)(80,30)
\qbezier(80,30)(90,20)(100,10)
\put(112,0){$x$}
\put(1,48){$f(x)$}
\put(37,32){\scriptsize $(2,3)$}
\put(55,43){\scriptsize $(3,4)$}
\put(81,32){\scriptsize $(4,3)$}
\put(98,4){\scriptsize $(5,1)$}
\end{picture}
\end{center}
